% ============================================================
% 第14章 高考综合大题训练（四道原创大题，按广东卷第17～20题的题型与分值设计）
% ============================================================
\chapter{高考综合大题训练}

本章的四道题都是原创题，分别练习实验探究、工艺流程、原理综合和有机综合，每题 14 分。建议学完全书后，每题用 15～20 分钟独立作答，再对照解析检查：知识是否会用，题给信息是否读全，结论是否有证据，计算是否符合守恒。第一次没有做完，可先补齐思路，再练限时作答。

\section{实验探究题}
\noindent\textbf{第1题}（14 分）某小组探究 \ce{FeCl3} 与 KI 的反应，并测定某 \ce{FeCl3} 溶液的浓度。

(1) 配制 100 mL 0.10 mol·L$^{-1}$ \ce{FeCl3} 溶液，需要的玻璃仪器除烧杯、玻璃棒、量筒外，还有\tk；配制时先用少量盐酸酸化的水溶解 \ce{FeCl3}，再转移并定容，酸化的目的是\tk。（2 分）

(2) 探究反应的限度：取 5 mL 0.10 mol·L$^{-1}$ KI 溶液，滴加 1 mL 0.10 mol·L$^{-1}$ \ce{FeCl3} 溶液，充分反应后分成两份。实验①：向第一份中加入 \ce{CCl4}，振荡静置，下层呈紫红色；实验②：向第二份中滴加 KSCN 溶液，溶液变红。
\quad ①写出 \ce{FeCl3} 与 KI 反应的离子方程式。（1 分）
\quad ②实验中为什么要让 KI 过量？实验②的现象说明了什么？（2 分）

(3) 探究浓度对平衡的影响：向实验②的红色溶液中加入少量 KI 固体，红色变浅。用平衡移动原理解释。（2 分）

(4) 测定浓度：取 25.00 mL 待测 \ce{FeCl3} 溶液于锥形瓶中，在适宜的弱酸性条件下加入过量 KI 溶液，充分反应后，用 0.1000 mol·L$^{-1}$ \ce{Na2S2O3} 标准溶液滴定（\ce{I2 + 2S2O3^2- = 2I- + S4O6^2-}），消耗 20.00 mL。滴定近终点时加入指示剂，忽略其他副反应。
\quad ①应选用的指示剂为\tk，滴定终点的现象为\tk。（2 分）
\quad ②该 \ce{FeCl3} 溶液的浓度为\tk mol·L$^{-1}$。（2 分）
\quad ③若放置过程中空气中的氧气氧化了部分过量 \ce{I-}，且忽略 \ce{I2} 挥发和其他副反应，测定结果\tk（填“偏高”“偏低”或“无影响”），原因是\tk。（2 分）

(5) 有同学认为，(4) 中“反应可逆”会使测定结果偏低。你是否同意？说明理由。（1 分）

\section{工艺流程题}
\noindent\textbf{第2题}（14 分）以含钴废料（主要含 \ce{Co2O3}、CoO，还含少量 \ce{Fe2O3}、\ce{Al2O3}、\ce{SiO2}）为原料制备草酸钴晶体（\ce{CoC2O4*2H2O}）和 \ce{Co3O4} 的流程如下：
\begin{center}
\begin{tikzpicture}[font=\footnotesize,node distance=0.45cm and 0.8cm,st/.style={draw,rounded corners=1.5pt,fill=lightblue,minimum height=0.6cm,inner sep=3pt},pr/.style={inner sep=1pt,align=center}]
\node[pr] (a) {含钴废料};
\node[st,right=of a] (b) {酸浸};
\node[st,right=of b] (c) {氧化};
\node[st,right=of c] (d) {调 pH};
\node[st,right=of d] (e) {沉钴};
\node[st,right=2.3cm of e] (f) {煅烧};
\node[pr,right=of f] (g) {\ce{Co3O4}};
\draw[-Stealth] (a)--(b); \draw[-Stealth] (b)--(c); \draw[-Stealth] (c)--(d); \draw[-Stealth] (d)--(e); \draw[-Stealth] (e)--node[below=1pt]{\ce{CoC2O4*2H2O}}(f); \draw[-Stealth] (f)--(g);
\node[pr,above=of b] (rb) {稀 \ce{H2SO4}\\\ce{Na2SO3}}; \draw[-Stealth] (rb)--(b);
\node[pr,above=of c] (rc) {\ce{NaClO3}}; \draw[-Stealth] (rc)--(c);
\node[pr,above=of d] (rd) {\ce{CoO}}; \draw[-Stealth] (rd)--(d);
\node[pr,above=of e] (re) {\ce{(NH4)2C2O4}}; \draw[-Stealth] (re)--(e);
\node[pr,below=of b] (r1) {滤渣1}; \draw[-Stealth] (b)--(r1);
\node[pr,below=of d] (r2) {滤渣2}; \draw[-Stealth] (d)--(r2);
\node[pr,below=of e] (r3) {滤液}; \draw[-Stealth] (e)--(r3);
\end{tikzpicture}
\end{center}
已知：25 ℃ 时，$K_\text{sp}[\ce{Fe(OH)3}]=2.8\times10^{-39}$，$K_\text{sp}[\ce{Al(OH)3}]=1.3\times10^{-33}$，$K_\text{sp}[\ce{Co(OH)2}]=1.6\times10^{-15}$；$\lg5.1\approx0.7$，$\lg1.26\approx0.1$。各沉淀计算忽略体积变化及配位等其他平衡；草酸钴晶体难溶于乙醇。

(1)“酸浸”时，提高浸取速率的措施有\tk（写一条）；滤渣1的主要成分是\tk。（2 分）

(2)“酸浸”时加入 \ce{Na2SO3} 的作用是\tk，写出 \ce{Co2O3} 参与反应的离子方程式。（2 分）

(3)“氧化”时 \ce{NaClO3} 将 \ce{Fe^2+} 氧化，写出离子方程式。已知 $K_\text{sp}[\ce{Fe(OH)2}]=4.9\times10^{-17}$，说明为什么要先把 \ce{Fe^2+} 氧化为 \ce{Fe^3+} 再调 pH。（2 分）

(4) 若“调 pH”前溶液中 $c(\ce{Co^2+})=0.10$ mol·L$^{-1}$，为使 \ce{Fe^3+}、\ce{Al^3+} 沉淀完全（离子浓度 $\le1.0\times10^{-5}$ mol·L$^{-1}$）而 \ce{Co^2+} 不沉淀，应控制 pH 的范围为\tk。（2 分）

(5)“沉钴”的离子方程式为\tk。用乙醇洗涤 \ce{CoC2O4*2H2O} 晶体的目的是\tk。（2 分）

(6) 取 18.3 g \ce{CoC2O4*2H2O}（$M=183$ g·mol$^{-1}$）在空气中加热，固体质量随温度变化：$300$ ℃ 时剩余 14.7 g，$450$ ℃ 以上剩余 8.03 g 并保持不变。
\quad ①300 ℃ 时剩余固体的化学式为\tk。（1 分）
\quad ②写出 300～450 ℃ 之间发生反应的化学方程式。（2 分）

(7) CoO 的晶胞与 NaCl 相同，晶胞边长为 0.426 nm，则 CoO 晶体的密度为\tk g·cm$^{-3}$（列出算式即可）。（1 分）

\section{原理综合题}
\noindent\textbf{第3题}（14 分）丙烷脱氢制丙烯是丙烯生产的重要途径，主反应和副反应如下：\\
主反应 Ⅰ：\ce{C3H8(g) <=> C3H6(g) + H2(g)} \quad $\Delta H_1$\\
副反应 Ⅱ：\ce{C3H8(g) <=> C2H4(g) + CH4(g)} \quad $\Delta H_2=+81$ kJ·mol$^{-1}$

(1) 已知键能（kJ·mol$^{-1}$）：C—C 348，C=C 615，C—H 413，H—H 436。估算 $\Delta H_1=$\tk kJ·mol$^{-1}$。（2 分）

(2) 本小问仅考虑反应Ⅰ，从纯丙烷出发。按 $\Delta H$、$\Delta S$ 判断，该反应在足够\tk（填“高温”或“低温”）时有利于自发进行；下列措施中，既能加快速率、又能提高丙烷平衡转化率的是\tk（填字母）。a．升高温度\quad b．恒温压缩容器\quad c．使用高效催化剂\quad d．保持相同温度、相同容积，增加丙烷初始投料量（2 分）

(3) 本小问忽略副反应，只将一定量纯 \ce{C3H8} 置于密闭体系，在恒温恒压条件下发生反应Ⅰ。下列能说明体系达到平衡状态的是\tk（填字母）。A．单位时间内生成 $n$ mol \ce{C3H6} 的同时生成 $n$ mol \ce{H2}\quad B．混合气体的密度不再变化\quad C．$n(\ce{C3H6}):n(\ce{H2})=1:1$\quad D．\ce{C3H8} 的体积分数不再变化（2 分）

(4) 某温度下，在总压恒为 100 kPa 的条件下，向反应器中通入 1 mol \ce{C3H8} 和 4 mol Ar，达到平衡时 \ce{C3H8} 的转化率为 50\%，丙烯的选择性为 80\%（丙烯的选择性 $=\dfrac{\text{转化为丙烯的 }\ce{C3H8}\text{ 的量}}{\text{转化的 }\ce{C3H8}\text{ 总量}}\times100\%$）。
\quad ①平衡时 \ce{H2} 的物质的量为\tk mol，气体总物质的量为\tk mol。（2 分）
\quad ②反应Ⅰ的平衡常数 $K_p=$\tk kPa（保留两位小数）。（2 分）

(5) 恒温、总压不变时，在含相同物质的量 \ce{C3H8} 的原料中增大 Ar 的比例，\ce{C3H8} 的平衡转化率增大；但 Ar 比例过大时，在相同催化剂用量下，\ce{C3H8} 的初始消耗速率明显下降。分别解释原因。（2 分）

(6) 长时间反应后催化剂表面会积炭而失活。研究发现，在原料气中加入适量 \ce{CO2} 可以减少积炭，用化学方程式表示其原理。（2 分）

\section{有机综合题}
\noindent\textbf{第4题}（14 分）对乙酰氨基酚（扑热息痛，\ce{HO-C6H4-NHCOCH3}，\ce{-OH} 与 \ce{-NHCOCH3} 处于苯环对位）是常用解热镇痛药，一种合成路线如下：
\tu{\begin{tikzpicture}[font=\small,>=Stealth]
\node[align=center] (ph) at (0,0) {苯酚};
\node[align=center] (a) at (4.2,0) {A：\ce{C6H5NO3}\\（分离出的对位产物）};
\node[align=center] (b) at (8.8,0) {B：\ce{C6H7NO}};
\node[align=center] (p) at (8.8,-1.7) {\ce{HO-C6H4-NHCOCH3}\\（两取代基在对位）};
\draw[->] (ph)--node[above]{① 稀 \ce{HNO3}}(a);
\draw[->] (a)--node[above]{② \ce{Fe}、\ce{HCl}}(b);
\draw[->] (b)--node[right]{③ \ce{(CH3CO)2O}}(p);
\end{tikzpicture}}{对乙酰氨基酚的合成路线；“对位”说明两个取代基在苯环上的相对位置}
路线中各步均经必要的分离提纯，步骤②酸性反应液经中和后得到 B。已知：\ce{-NH2} 与乙酸酐的反应活性远大于酚羟基；\ce{R-NO2 ->[\ce{Fe},\ \ce{HCl}] R-NH2}；酰胺键 \(\mathrm{RCONHR'}\) 在足量 NaOH 溶液中共热水解，生成 \ce{RCOONa} 和 \(\mathrm{R'NH_2}\)。

(1) 苯酚与稀硝酸反应生成 A 的反应类型为\tk；A 的名称为\tk。（2 分）

(2) 反应②中有机物发生了\tk（填“氧化”或“还原”）反应；B 的结构简式为\tk。（2 分）

(3) 写出反应③的化学方程式。（2 分）

(4) 对乙酰氨基酚分子中碳原子的杂化方式有\tk；1 mol 对乙酰氨基酚与足量 NaOH 溶液共热，最多消耗 NaOH\tk mol。（2 分）

(5) 对乙酰氨基酚的同分异构体中，同时满足下列条件的有\tk{}种。①含苯环，苯环上只有两个取代基；②其中一个取代基为 \ce{-NH2}；③能发生银镜反应，也能在 NaOH 溶液中水解。（2 分）

(6) 以苯胺为原料在苯环对位引入硝基等基团时，通常先用乙酸酐把 \ce{-NH2} 转化为 \ce{-NHCOCH3}（乙酰苯胺），反应结束后再水解恢复 \ce{-NH2}。说明这样做的目的。（2 分）

(7) 乙酸酐可由乙酸制得。以乙烯为原料（无机试剂任选）制备乙酸，合成路线为：乙烯 $\to$ \tk $\to$ \tk $\to$ 乙酸（写出中间产物的结构简式并注明各步反应条件）。（2 分）

\zhangmo{参考答案与评分要点}
\begin{jiexilist}
\item 第1题\\
(1) 100 mL 容量瓶、胶头滴管（1 分）；抑制 \ce{Fe^3+} 水解（1 分）。\\
(2) ①\ce{2Fe^3+ + 2I- <=> 2Fe^2+ + I2}（1 分，注意可逆号）。②KI 过量，若反应能进行到底，不足量的 \ce{Fe^3+} 应完全反应（1 分）；已有 \ce{I2} 生成，且加 KSCN 变红说明 \ce{Fe^3+} 仍有剩余，支持该反应在此条件下存在限度（1 分）。\\
(3) 加入 KI，$c(\ce{I-})$ 增大，\ce{2Fe^3+ + 2I- <=> 2Fe^2+ + I2} 正向移动，$c(\ce{Fe^3+})$ 减小；有色配合物的形成平衡随之逆向移动，红色变浅（2 分，两个平衡各 1 分）。可把后一平衡简化写为 \ce{Fe^3+ + SCN- <=> [FeSCN]^2+}，实际体系中可同时存在不同配位数的有色微粒。\\
(4) ①淀粉溶液（1 分）；滴入最后半滴 \ce{Na2S2O3} 溶液时，溶液蓝色褪去，且半分钟内不恢复（1 分）。②关系式 \ce{2Fe^3+}～\ce{I2}～\ce{2S2O3^2-}，$n(\ce{Fe^3+})=n(\ce{S2O3^2-})=2.000\times10^{-3}$ mol，$c=\dfrac{2.000\times10^{-3}}{0.02500}=0.08000$ mol·L$^{-1}$（2 分）。③偏高（1 分）；空气中的 \ce{O2} 将 \ce{I-} 氧化为 \ce{I2}（\ce{4I- + O2 + 4H+ = 2I2 + 2H2O}），消耗 \ce{Na2S2O3} 偏多（1 分）。\\
(5) 不同意（1 分，理由合理即可）：KI 大量过量，且滴定过程中 \ce{I2} 不断被 \ce{S2O3^2-} 消耗，使平衡持续正向移动，\ce{Fe^3+} 最终能完全反应。

\item 第2题\\
(1) 粉碎废料、适当升温、适当增大硫酸浓度、搅拌（任一条，1 分）；\ce{SiO2}（1 分）。\\
(2) 作还原剂，将 $+3$ 价 Co 还原为 \ce{Co^2+}（1 分）；\ce{Co2O3 + SO3^2- + 4H+ = 2Co^2+ + SO4^2- + 2H2O}（1 分）。（\ce{Fe^3+} 同时被还原为 \ce{Fe^2+}。）\\
(3) \ce{ClO3^- + 6Fe^2+ + 6H+ = Cl- + 6Fe^3+ + 3H2O}（1 分）；\ce{Fe(OH)2} 与 \ce{Co(OH)2} 的 $K_\text{sp}$ 接近，\ce{Fe^2+} 沉淀完全时 \ce{Co^2+} 也已大量沉淀，无法通过调 pH 分离；\ce{Fe(OH)3} 的 $K_\text{sp}$ 极小，\ce{Fe^3+} 在 pH 约 2.8 时即沉淀完全，可与 \ce{Co^2+} 分离（1 分）。\\
(4) $4.7\le$pH$<7.1$（2 分）。\ce{Fe(OH)3} 的 $K_\text{sp}$ 更小，\ce{Al^3+} 完全沉淀时 \ce{Fe^3+} 早已沉淀完全：$c(\ce{OH-})=\sqrt[3]{\dfrac{1.3\times10^{-33}}{1.0\times10^{-5}}}\approx5.1\times10^{-10}$，pOH$\approx9.3$，pH$\approx4.7$；\ce{Co^2+} 开始沉淀时 $c(\ce{OH-})=\sqrt{\dfrac{1.6\times10^{-15}}{0.10}}\approx1.26\times10^{-7}$，pOH$\approx6.9$，pH$\approx7.1$。\\
(5) \ce{Co^2+ + C2O4^2- + 2H2O = CoC2O4*2H2O v}（1 分）；减少晶体的溶解损失，乙醇易挥发，便于晶体干燥（1 分）。\\
(6) ①$n=0.1$ mol，300 ℃ 时固体 $M=147$ g·mol$^{-1}$，为 \ce{CoC2O4}（失去全部结晶水）（1 分）。②8.03 g 中含 Co 0.1 mol（5.9 g），O 为 $8.03-5.9=2.13$ g，约 0.133 mol，$n(\ce{Co}):n(\ce{O})=3:4$，为 \ce{Co3O4}：\ce{3CoC2O4 + 2O2 ->[\text{△}] Co3O4 + 6CO2}（2 分）。\\
(7) $\rho=\dfrac{4\times75}{N_\text{A}\times(0.426\times10^{-7})^3}$（约 6.4 g·cm$^{-3}$）（1 分）。

\item 第3题\\
(1) $+123$（2 分）。断键：1 个 C—C 和 2 个 C—H（反应物中多出的部分）；成键：1 个 C=C 和 1 个 H—H。$\Delta H_1=348+2\times413-615-436=+123$ kJ·mol$^{-1}$。（可只计算发生变化的键。）\\
(2) 高温（$\Delta H>0$、$\Delta S>0$，温度足够高时 $\Delta G=\Delta H-T\Delta S<0$）（1 分）；a（1 分）。b 使平衡逆移；c 只加快达到平衡，不改变平衡组成。对 d，设初始浓度为 $c_0$、平衡转化率为 $\alpha$，则 $K_c=c_0\alpha^2/(1-\alpha)$。恒温时 $K_c$ 不变，增大 $c_0$ 必使 $\alpha$ 减小；这与加入反应物后平衡可正向移动并不矛盾：绝对转化量和转化率的分母同时改变，不能只用移动方向判断转化率。\\
(3) BD（2 分，选对一个得 1 分，错选不得分）。本小问只有反应Ⅰ。设初始有 1 mol 丙烷、转化 $x$ mol，则气体总量为 $1+x$ mol；恒温恒压下 $V\propto1+x$，总质量不变，故密度随 $x$ 增大而减小。密度稳定说明转化程度不再变化，B 正确。丙烷体积分数为 $(1-x)/(1+x)$，也随 $x$ 改变，D 正确。A 中两个“生成”都是正反应方向；C 中两种产物始终按 $1:1$ 形成，均不能说明正、逆反应速率相等。若同时存在多个独立反应，一个总量不变未必足以证明整个体系平衡。\\
(4) ①转化 \ce{C3H8} 0.5 mol，其中 0.4 mol 生成丙烯，0.1 mol 发生副反应。\ce{H2} 0.4 mol（1 分）；平衡时 \ce{C3H8} 0.5、\ce{C3H6} 0.4、\ce{H2} 0.4、\ce{C2H4} 0.1、\ce{CH4} 0.1、Ar 4，总 5.5 mol（1 分）。②$K_p=\dfrac{(\frac{0.4}{5.5}\times100)^2}{\frac{0.5}{5.5}\times100}=\dfrac{0.16\times100}{0.5\times5.5}\approx5.82$ kPa（2 分）。\\
(5) 总压不变时增大 Ar 比例，反应物、生成物分压同比减小，相当于减压，平衡向气体分子数增大的正方向移动，转化率增大（1 分）；Ar 比例过大，\ce{C3H8} 分压（浓度）过低，反应速率下降（1 分）。\\
(6) \ce{C + CO2 ->[\text{高温}] 2CO}（2 分）。

\item 第4题\\
(1) 取代反应（硝化反应）（1 分）；对硝基苯酚（4-硝基苯酚）（1 分）。\\
(2) 还原（1 分）；\ce{HO-C6H4-NH2}（对氨基苯酚，对位）（1 分）。\\
(3) \ce{HO-C6H4-NH2 + (CH3CO)2O -> HO-C6H4-NHCOCH3 + CH3COOH}（2 分，产物、配平各 1 分）。\\
(4) $\mathrm{sp}^2$、$\mathrm{sp}^3$（1 分）；2（1 分）。酚羟基与 NaOH 中和消耗 1 mol；酰胺基在碱性条件下水解生成 \ce{CH3COONa}（与对氨基苯酚钠），消耗 1 mol，共 2 mol。\\
(5) 3（2 分）。先由结构得分子式 \ce{C8H9NO2}。苯环上两个位置已被取代，环的组成为 \ce{C6H4}；扣除 \ce{-NH2} 后，另一取代基只能为 \ce{C2H3O2}。两个 O 既要提供能发生碱性水解的酯结构，又要满足银镜反应，需形成甲酸酯；从苯环向外写，这个取代基是 \ce{-CH2-O-C(=O)H}，即 \ce{HCOOCH2-}。它与 \ce{-NH2} 有邻、间、对 3 种位置关系。不能把 \ce{-CH2-O-C(=O)H} 误画成含 O—O 键的结构。\\
(6) \ce{-NH2} 还原性强，直接硝化时会被浓硝酸氧化（1 分）；乙酰化后 \ce{-NHCOCH3} 不易被氧化，起“保护氨基”的作用，同时它对苯环的活化作用减弱、体积较大，可减少多取代并有利于得到对位产物，反应后水解即可恢复 \ce{-NH2}（1 分）。\\
(7) \ce{CH2=CH2 ->[H2O][\text{催化剂，加热、加压}] CH3CH2OH ->[O2, Cu][\text{△}] CH3CHO ->[O2][\text{催化剂，△}] CH3COOH}（中间产物各 1 分；条件缺失酌情扣分）。
\end{jiexilist}
